% ============================================================================
% Hindi UI and structural strings for the One Chemistry Book series.
% Loaded by styles/onechemistry.sty when \booklang is hi.
% UTF-8 Devanagari throughout. Build with XeLaTeX (see *_hi.tex entries).
% ============================================================================

% Theorem / statement environment titles
\newcommand{\omnameDefinition}{परिभाषा}
\newcommand{\omnameTheorem}{प्रमेय}
\newcommand{\omnameProposition}{प्रतिज्ञप्ति}
\newcommand{\omnameLemma}{प्रमेयिका}
\newcommand{\omnameCorollary}{उपप्रमेय}
\newcommand{\omnameMethod}{विधि}
\newcommand{\omnameExample}{उदाहरण}
\newcommand{\omnameNotation}{संकेतन}
\newcommand{\omnameRemark}{टिप्पणी}
\newcommand{\omnameExercise}{अभ्यास}
\newcommand{\omnameExercises}{अभ्यास}
\newcommand{\omnameProblem}{समस्या}
\newcommand{\omnameProblems}{समस्याएँ}
\newcommand{\omnameProof}{उपपत्ति}

% Admitted results and solutions appendix headers
\newcommand{\omadmittedtext}{इस स्तर पर स्वीकृत।}
\newcommand{\omsolutionof}{हल ---}
\newcommand{\ompage}{पृष्ठ}
\newcommand{\omnameGotoSolution}{हल पृ.}

% Misc text macros
\newcommand{\st}{\text{ ऐसा कि }}

% Cover / title page
\newcommand{\bookmaintitle}{One Chemistry Book}
\newcommand{\booktagline}{सब पर राज करने वाली एक रसायन विज्ञान की पुस्तक}
\newcommand{\bookdescription}{विद्यालय की पहली कक्षा से स्नातक की डिग्री के अंत तक रसायन विज्ञान}
\newcommand{\bookcontributors}{One Chemistry Book के योगदानकर्ता}
\newcommand{\bookversionlabel}{संस्करण}

% Colophon
\newcommand{\omcolophonsource}{स्रोत कोड, समस्याएँ और योगदान:}
\newcommand{\omcolophonlicense}{जहाँ अन्यथा उल्लेख न हो, यह पुस्तक Creative Commons CC BY-NC-SA 4.0 लाइसेंस के अंतर्गत उपलब्ध है:}
\newcommand{\omcolophoncredit}{कृपया इस प्रकार श्रेय दें: One Chemistry Book, One Course (one-course.com)।}

% hi Book 1 agent, 2026-10-04: the recall title is pronoun-free ("already
% known") because the same box opens chapters written for a child (tum
% imperatives, grades 1-9) and for an upper-secondary reader (aap, grades
% 10-12); "jo aap pahle se jaante hain" fitted only the second.
% Chemistry boxes (styles/onechemistry.sty): the recall box that opens a
% chapter building on another one, the lab box, the history box, the safety
% box. Placeholder wording by the coordinator (2026-10-04); the Book 1
% `hi` agent owns it and settles the register (the recall title addresses
% the reader in the same register as the book's exercise stems).
\newcommand{\omnameRecall}{पहले से ज्ञात}
\newcommand{\omnameInTheLab}{प्रयोगशाला में}
\newcommand{\omnameHistory}{इतिहास}
\newcommand{\omnameSafety}{सुरक्षा}

% Solutions appendix part title
\newcommand{\omsolutionspart}{अभ्यासों के हल}

% Part titles (Grades 1–12; parallel to English Grade N)
\@namedef{omstr@part.grade1}{कक्षा 1 (आयु 6--7)}
\@namedef{omstr@part.grade2}{कक्षा 2 (आयु 7--8)}
\@namedef{omstr@part.grade3}{कक्षा 3 (आयु 8--9)}
\@namedef{omstr@part.grade4}{कक्षा 4 (आयु 9--10)}
\@namedef{omstr@part.grade5}{कक्षा 5 (आयु 10--11)}
\@namedef{omstr@part.grade6}{कक्षा 6 (आयु 11--12)}
\@namedef{omstr@part.grade7}{कक्षा 7 (आयु 12--13)}
\@namedef{omstr@part.grade8}{कक्षा 8 (आयु 13--14)}
\@namedef{omstr@part.grade9}{कक्षा 9 (आयु 14--15)}
\@namedef{omstr@part.grade10}{कक्षा 10 (आयु 15--16)}
\@namedef{omstr@part.grade11}{कक्षा 11 (आयु 16--17)}
\@namedef{omstr@part.grade12}{कक्षा 12 (आयु 17--18)}

% Part titles (Bachelor Years 1–3)
\@namedef{omstr@part.bachelor1}{स्नातक वर्ष 1 (आयु 18--19)}
\@namedef{omstr@part.bachelor2}{स्नातक वर्ष 2 (आयु 19--20)}
\@namedef{omstr@part.bachelor3}{स्नातक वर्ष 3 (आयु 20--21)}

% Plural names + \omnameChapter, used by the \crefname declarations in
% onechemistry.sty.
\newcommand{\omnameDefinitions}{परिभाषाएँ}
\newcommand{\omnameTheorems}{प्रमेय}
\newcommand{\omnamePropositions}{प्रतिज्ञप्तियाँ}
\newcommand{\omnameLemmas}{प्रमेयिकाएँ}
\newcommand{\omnameCorollarys}{उपप्रमेय}
\newcommand{\omnameMethods}{विधियाँ}
\newcommand{\omnameExamples}{उदाहरण}
\newcommand{\omnameNotations}{संकेतन}
\newcommand{\omnameRemarks}{टिप्पणियाँ}
\newcommand{\omnameChapter}{अध्याय}
\newcommand{\omnameChapters}{अध्याय}
% Section, for \cref{sec:…} (2026-10-04: it had no \crefname, so cleveref
% printed its English built-in "Section" in every edition).
\newcommand{\omnameSection}{अनुभाग}
\newcommand{\omnameSections}{अनुभाग}

% siunitx and cleveref keep their English conjunctions unless told otherwise,
% and Hindi has no babel here to localise them (see the note below). So
% \qtyrange printed "20 to 480 Hz" and \cref{a,b} printed "अध्याय 25 and 29"
% in the middle of Devanagari prose -- 19 such sites in the shipped Book 2
% Hindi edition, invisible to every structural gate because they are produced
% by the macros, not written in the sources.
% siunitx typesets range phrases and list separators in MATH MODE, where a
% bare string loses its spaces (\qtyrange inside $...$ printed "1a100") and a
% non-ASCII character is worse: "a" expands to \`a, which LaTeX reports as
% "Command \` invalid in math mode" and drops from the page. siunitx's own
% defaults are wrapped in \text{} for exactly this reason; these overrides
% must be too. Found on Biology Book 3 fr, 2026-09-05, where the book's 56
% \qtyrange sites made it visible; it was latent in every language edition of
% every book in this repo.
\sisetup{range-phrase={\text{ से }}, list-pair-separator={\text{ और }},
         list-final-separator={\text{ और }}}
% Deferred to the END of begindocument: cleveref creates these macros inside
% its own \AtBeginDocument, so a \renewcommand at package-load time fails with
% "Command \crefpairconjunction undefined", and a plain \AtBeginDocument only
% wins if ours was registered after cleveref's. Same technique as the Arabic
% block in onechemistry.sty, for the same reason.
\AddToHook{begindocument/end}{%
  \renewcommand{\crefpairconjunction}{ और }%
  \renewcommand{\creflastconjunction}{ और }%
  \renewcommand{\crefpairgroupconjunction}{ और }%
  \renewcommand{\creflastgroupconjunction}{ और }%
  \renewcommand{\crefrangeconjunction}{ से }%
}

% LaTeX's own structural names (babel is not used for Hindi).
\renewcommand{\chaptername}{अध्याय}
\renewcommand{\partname}{भाग}
\renewcommand{\contentsname}{विषय-सूची}
\renewcommand{\appendixname}{परिशिष्ट}
\renewcommand{\indexname}{अनुक्रमणिका}

% Periodic-table legend (\omperiodictable[legend], styles/onechemistry.sty).
% Coordinator wording, 2026-10-04 (it used to be hard-coded English in the
% style file, so no edition could translate it); the Book 1 agent of this
% language checks it against the book's own terminology.
\newcommand{\omnamePTmetal}{धातुएँ}
\newcommand{\omnamePTmetalloid}{उपधातुएँ}
\newcommand{\omnamePTnonmetal}{अधातुएँ}
\newcommand{\omnamePTs}{s-ब्लॉक}
\newcommand{\omnamePTp}{p-ब्लॉक}
\newcommand{\omnamePTd}{d-ब्लॉक}
\newcommand{\omnamePTf}{f-ब्लॉक}
\newcommand{\omnamePTalkali}{क्षार धातुएँ}
\newcommand{\omnamePTalkaline}{क्षारीय मृदा धातुएँ}
\newcommand{\omnamePThalogen}{हैलोजन}
\newcommand{\omnamePTnoble}{उत्कृष्ट गैसें}
